% =====================================================================
%  GMN3 §3  希尔伯特空间、halo 与穿墙（Hilbert spaces, halos, and wall-crossing）
%  §3.1–3.3；原文 Lamination-Paper.tex 1029--1506 行；PDF p.14--21（印刷 p.13--20）
%  编号公式 (3.1)--(3.24) 共 24 条；图 1 幅（spectrum）；脚注 9 个
%  跨块引用（式 (2.10)、(10.33)，第 4 节，附录 D）暂写原文编号并加 % TODO，
%  合稿时统一改 \eqref / \ref；label 名一律照抄原文。
%  ！注意：(3.6) 沿用原文 label \label{eq:bps-bound}，与 ch01 给 (1.2)
%  自拟的同名 label 撞名，合稿时由 Lead 择一改名（本块遵从原文）。
% =====================================================================

\section{希尔伯特空间、halo 与穿墙}
\label{sec:Hilbert-Halo-WC}

\subsection{带框 BPS 态}
\label{sec:framed-BPS-states}

此前我们只讨论了 $\CN=2$ 理论的紫外方面。现在我们要过渡到红外。先交代一点记号。我们把库仑分支记作 $\CB$，把它的奇异除子记作 $\CB^{sing}$，并把「原味」电磁荷与味荷构成的格记作 $\Gamma$。\footnote{众所周知，若沿 $\CB^*:=\CB-\CB^{sing}$ 中的非平凡路径对 $u$ 作解析延拓，$\Gamma$ 会经历单值变换。因此严格说来，该谈论的是 $\Gamma_u$。用数学的话说，$\Gamma$ 是 $\CB^*$ 上的一个格丛，并定义了一个「局部系统」（local system），也就是说，它带有一个平坦联络。人们可以在 $\CB-\CB^{sing}$ 的万有覆盖 $\widehat\CB$ 上工作，在那里 $\Gamma$ 可以被平凡化；也可以直接在底空间上工作，只要记住存在非平凡的单值性。这是个人口味问题，我们留给读者。更多的背景见 \cite{Gaiotto:2008cd,Gaiotto:2009hg} 及其中的文献。}

现在我们来研究理论在 $\IR^{1,3}$ 上的希尔伯特空间：无穷远处的真空由 $u\in\CB$ 标定，并且在原点 $x^i=0$ 处插入了一个简单线算子 $L_\zeta$。这个希尔伯特空间我们记作 $\CH_{u,L,\zeta}$，它是 $osp(4^*\vert 2)_\zeta$ 的庞加莱子超代数的一个表示。此外，它按电荷分次：
\begin{equation}\label{eq:Grade-Hilbert}
\CH_{u,L,\zeta} = \bigoplus_{\gamma \in \Gamma_L}
\CH_{u,L,\zeta,\gamma}.
\end{equation}
这里 $\Gamma_L$ 是 $L$ 存在时的电磁荷与味荷的泊松格，它是「原味」电荷之格 $\Gamma$ 的一个挠子。也就是说，它具有如下形式
\begin{equation}\label{eq:Lattice-Shift}
\Gamma_L = \Gamma + \gamma_L
\end{equation}
其中 $\gamma_L \in \Gamma \otimes \IR$，且对一切 $\gamma \in \Gamma$ 有 $\langle \gamma_L, \gamma\rangle \in \IZ$。若用路径积分表述来说，\eqref{eq:Grade-Hilbert} 中对 $\Gamma$ 的直和就变成对拓扑扇区的求和。

\eqref{eq:Grade-Hilbert} 中对 $\gamma$ 求和的必要性，在弱耦合的纯 $SU(N)$ 理论中很容易看出来：取库仑分支上足够远的位置，使规范对称性被强破缺为 $U(1)^{N-1}$。例如考虑 Wilson 线算子 (2.10)% TODO 合稿改 \eqref{eq:Susy-Wilson-Zeta}。在朴素的经典极限下，这个算子的 vev 将是对红外理论中若干超对称 Wilson 线所对应的 vev 之和，这些 Wilson 线由表示 $\CR$ 的权来标记。\footnote{也许出人意料的是，在具体例子中我们会看到这个朴素的答案漏掉了一些东西：对权求和时还有额外的贡献。一个简单的例子是纯 $SU(2)$ 规范理论中基础表示下的 Wilson 算子，其答案我们在下面的式 (10.33)% TODO 合稿改 \eqref{eq:wilson-vev-nf0} 中给出。对应于基础表示之权的项有两项，但还有第三项。我们会把这第三项归因于粒子与线算子形成的一个有趣的束缚态。} 这个例子同时也说清了 \eqref{eq:Lattice-Shift} 中那个平移的物理起源。考虑纯 $SU(2)$ 规范理论中基础表示 Wilson 线算子的例子：所有「原味」BPS 态都带偶数电荷（在使 W 玻色子带电荷 $2$ 的约定下），但在一个基础表示 Wilson 圈存在时，一个态应当能够带奇数电荷。往系统里添加「原味」BPS 粒子的自由，使「$\Gamma_L$ 应当是 $\Gamma$ 的挠子」这一点变得显而易见。

线算子上的相互局域条件表明：对一组合乎一致的线算子 $L\in\CL$，相应的 $\Gamma_L$ 应全部落在同一个格 $\Gamma_{\CL}$ 之中，使得
\begin{equation}\label{eq:Consist-LO-Latt}
\Gamma \subset \Gamma_{\CL},
\end{equation}
并且 $\Gamma_{\CL}$ 上有一个整值反对称形式，它限制到 $\Gamma$ 上就是标准的那一个。

扇区 $\CH_{u,L,\zeta,\gamma}$ 中的每一个都是 $osp(4^*\vert 2)_\zeta$ 的庞加莱子超代数的一个表示。当我们以 $x^0$ 为常数的时片来量子化时，超对称算子满足厄米性条件
\begin{equation}
\begin{split}
(\CR_{1}^{~1})^\dagger & = - \CR_{2}^{~2}, \\
(\CR_{1}^{~2})^\dagger & =   \CR_{2}^{~1}.
\end{split}
\end{equation}
此外，在从超共形点形变而来的理论中，希尔伯特空间会有一些扇区带有 $\CN=2$ 庞加莱超对称的非零中心荷。算子 $\CR_{\alpha}^{~A}$ 满足代数
\begin{equation}
\{ \CR_{\alpha}^{~A} , \CR_{\beta}^{~B} \} = 4 \left( E + {\rm Re}(Z/\zeta) \right) \epsilon_{\alpha\beta} \epsilon^{AB}
\end{equation}
其中 $E=P^0$ 是能量算子。把它与厄米性条件结合起来，就得到 BPS 界
\begin{equation} \label{eq:bps-bound-3d}
E + {\rm Re}\, (Z / \zeta) \geq 0,
\end{equation}
其中 $Z = Z_\gamma(u)$ 是与真空 $u\in \CB$、扇区 $\gamma$ 相联系的标准中心荷。

现在可以给出本文的主要定义之一了：我们定义\emph{带框 BPS 态}（framed BPS state）为使这个 BPS 界取到等号的态，即 $E = -{\rm Re}\, (Z / \zeta)$。（这个名字的动机来自带框箭图表示（framed quiver representation）的概念。）我们可以把 \eqref{eq:Grade-Hilbert} 特殊到由带框 BPS 态构成的子空间：
\begin{equation}\label{eq:Grade-BPS-Hilbert}
\CH_{u,L,\zeta}^{\rm BPS} = \bigoplus_{\gamma \in \Gamma_L}
\CH_{u,L,\zeta,\gamma}^{\rm BPS}.
\end{equation}

注意 \eqref{eq:bps-bound-3d} 与标准的 BPS 界 $E \geq \vert Z \vert$ 不同。理解 \eqref{eq:bps-bound-3d} 有一个很好的启发式办法：它用的是「把 Wilson–'t Hooft 算子看作插入一个无穷重双荷子」这一标准解释的红外版本。把格 $\Gamma$ 用一个额外的味荷 $\gamma_f$ 扩张，并考虑一个非常重的粒子，其电荷为 $\gamma_f - \gamma$、中心荷为 $Z = \zeta M  -  Z_\gamma$，其中 $M>0$。在 $M \to +\infty$ 极限下，重整化后的 BPS 界正好是
\begin{equation}\label{eq:heuristic-bound}
E \geq  \lim_{M\to +\infty} \left(\vert \zeta M  -  Z_\gamma\vert - M \right)= - {\rm Re}(Z_\gamma/\zeta).
\end{equation}

总电荷为 $\gamma$ 的 BPS 粒子，其能量有下界 $\vert Z_\gamma(u)\vert$。然而我们在下面会论证：在线算子存在时，可以有一些有趣的量子态，其能级低于 $\vert Z_\gamma(u)\vert$、但高于界 $-{\rm Re}\, (Z / \zeta)$。从远处看，它们像是位于线算子所在位置、电荷为 $\gamma_c$ 的「芯粒子」（core particle），与一个（或多个）电荷为 $\gamma_h$ 的「halo 粒子」相互作用，且 $\gamma_c + \gamma_h = \gamma$。在那些 halo 粒子并未束缚的情形——也就是说，其波函数并未本质上局限在空间中一个有限区域内——能量从 $-{\rm Re}\, (Z_{\gamma_c}(u) / \zeta)+ \vert Z_{\gamma_h}(u)\vert$ 开始，构成连续谱。此外还有束缚态，它们类比于超引力的多中心束缚态 \cite{Denef:2000nb,Denef:2000ar}。对这些态而言，halo 粒子的波函数本质上局限在空间中一个紧致区域内。这类态具有离散的能谱。只要不存在使界 $-{\rm Re}\, (Z_{\gamma_h}(u) / \zeta) = \vert Z_{\gamma_h}(u)\vert$ 取到等号的 halo 电荷，带框 BPS 态与连续谱之间就存在非零能隙。图 \ref{fig:spectrum} 示意性地画出了能谱。

\insfig{spectrum}{在插入了线算子 $L_\zeta$ 的理论中，电荷为 $\gamma$ 的诸态之能量谱的示意图。最下方是带框 BPS 态；接着是束缚在线算子上的离散激发态；最后是各种非束缚态的连续谱，它们对应于把 $\gamma$ 分解为 $\gamma_c + \gamma_h$ 的不同方式。带框 BPS 态与各连续谱之间有安全的间隔，除非存在某个 $\gamma_h$ 使 $-{\rm Re}\, (Z_{\gamma_h}(u) / \zeta)= \vert Z_{\gamma_h}(u)\vert$。}

条件 $-{\rm Re}\, (Z_\gamma(u) / \zeta)= \vert Z_{\gamma}(u)\vert$ 定义了\emph{BPS 墙}：\footnote{在 \cite{Gaiotto:2008cd,Gaiotto:2009hg} 中我们考虑的是「BPS 射线」，它们不过是固定 $u$ 时 BPS 墙到 $\zeta$-球面的投影。}
\begin{equation}\label{eq:Def-BPS-Walls}
\widehat W(\gamma ):=\{ (u,\zeta) \vert Z_{\gamma}(u)/\zeta \in
\IR_- \}\subset \widetilde \CB \times \widetilde\IC^*.
\end{equation}
当我们在空间 $\tilde \CB \times \IC^*$ 中变动参数时，一旦击中某道 BPS 墙，带框 BPS 态就会与连续谱混合。

\subsection{带框 BPS 简并度与保护自旋特征标}
\label{subsec:BPS-Deg-PSC}

带框 BPS 态按空间转动李代数 $so(3)$ 的表示变换。在固定分次 $\gamma$ 下它们是有限维的，于是我们可以定义一个简单的指标，来计数带框 BPS 简并度：
\begin{equation}
\fro(u,L_\zeta, \gamma) := {\rm Tr}_{\CH_{u,L,\zeta,\gamma}^{\rm
BPS}} (-1)^{(2J_3)}
\end{equation}
事实上，我们还可以定义一个更精细的量来计数带框 BPS 态。我们希望定义一个能记录自旋信息的量。简单的自旋特征标 ${\rm Tr}\, y^{2J_3}$ 朴素地看就够用了，但若若干 BPS 态配对成超对称代数的一个非保护表示，它就可能发生难以预测的跳变。为得到一个更稳健的指标，我们回想：带框 BPS 态也按 $su(2)_R$ 变换，而四个超荷 $\CR_\alpha^A$ 中有一个，在自旋 $so(3)$ 与 $su(2)_R$ 的对角组合下是单态，即
\begin{equation}
\CQ := \epsilon_{\alpha A} \CR_{\alpha}^A.
\end{equation}
于是我们可以定义所谓的\emph{（带框）保护自旋特征标}，简称 PSC：\footnote{我们感谢 Juan Maldacena 建议了这个改进的指标。}
\begin{equation}\label{eq:Def-Framed-PSC}
\fro(u,L_\zeta, \gamma;y) := {\rm Tr}_{\CH_{u,L,\zeta,\gamma}^{\rm
BPS}} (-1)^{2J_3} (-y)^{2\CJ_3} = {\rm Tr}_{\CH_{u,L,\zeta,\gamma}^{\rm
BPS}} y^{2J_3} (-y)^{2I_3}
\end{equation}
其中 $I_3$ 是一个 R 对称生成元，且
\begin{equation}
\CJ_3 := J_3 + I_3.
\end{equation}
PSC 是一个指标，因为 $\CQ$ 在 $\CJ_a$ 下是单态、与 $(-1)^F$ 反对易，并且在 $\CR_{\alpha}^A$ 代数的长表示上可逆。通常我们会把记号缩短，把对 $u$ 的依赖略去不写。

特殊到 $y=-1$，我们就回到带框 BPS 简并度 $\fro(u,L_\zeta, \gamma)$。由于超对称小表示的刚性，只要没有态进入或离开希尔伯特空间，它们就在参数的形变下不变。不过，仍会出现穿墙现象，类比于没有线算子时 BPS 态的通常穿墙。我们在第 \ref{subsec:Halos} 节中研究这一点。

对通常的「原味」BPS 态也可以定义 PSC。「原味」PSC 的定义基于这样一个论断：量
\begin{equation}\label{eq:PSC-one}
{\rm Tr}_{\CH}(2 J_3) (-1)^{2J_3} (-y)^{2\CJ_3}
\end{equation}
在 $\CN=2$ 代数的长表示 $\CH$ 上为零。为证明这一点，考虑 $so(3)\oplus su(2)_R$ 的特征标 ${\rm Tr}\, x_1^{2J_3} x_2^{2I_3}$。由它出发，通过
\begin{equation}\label{eq:Inermed-PSC}
{\rm Tr}(2 J_3) (-1)^{2J_3} (-y)^{2\CJ_3} = x_1\frac{\partial}{\partial x_1}\left({\rm Tr}\, x_1^{2J_3} x_2^{2I_3}\right)\vert_{x_1=-x_2 = y}
\end{equation}
就得到 \eqref{eq:PSC-one}。现在，在有质量粒子的小超代数的表示论中，长表示作为 $so(3)\oplus su(2)_R$ 的表示具有如下形式
\begin{equation}
\rho_{hh}\otimes \rho_{hh}\otimes \mathfrak{h}_{\ell}
\end{equation}
其中 $\mathfrak{h}_{\ell}$ 是 $so(3)\oplus su(2)_R$ 的任意有限维表示，而 $\rho_{hh}$ 是\emph{半超多重态}表示 $\rho_{hh}\cong (\half;0) \oplus (0;\half)$。另一方面，小超代数的短表示作为 $so(3)\oplus su(2)_R$ 的表示具有如下形式
\begin{equation}
\rho_{hh} \otimes \mathfrak{h}_{s}
\end{equation}
其中 $\mathfrak{h}_{s}$ 是 $so(3)\oplus su(2)_R$ 的任意有限维表示。现在注意
\begin{equation}
{\rm Tr}_{\rho_{hh}}  x_1^{2J_3} x_2^{2I_3} = x_1 + x_1^{-1} + x_2 + x_2^{-1}
\end{equation}
因此 \eqref{eq:Inermed-PSC} 在长表示上为零，但在短表示上非零。事实上，它在短表示上的值恰好是
\begin{equation}
(y- y^{-1}) {\rm Tr}_{\mathfrak{h}_s} y^{2J_3} (-y)^{2I_3} .
\end{equation}
于是，我们可以通过
\begin{equation}
(y- y^{-1}) \Omega(u,\gamma;y) := {\rm Tr}(2 J_3) (-1)^F (-y)^{2\CJ_3}.
\end{equation}
来定义一个（不带框的）保护自旋特征标 $\Omega(u,\gamma;y)$，即
\begin{equation}
\Omega(u,\gamma;y) = {\rm Tr}_{\mathfrak{h}_s} y^{2J_3} (-y)^{2I_3} .
\end{equation}
注意，若 $\mathfrak{h}_s$ 作为 $su(2)_R$ 表示的同型分解只含单态，则保护自旋特征标\emph{等同于} $\mathfrak{h}_s$ 的自旋特征标。于是我们立刻看出：一个标准超多重态有 $\mathfrak{h}_s = (0;0)$，故其 PSC 为 $\Omega(u,\gamma;y)=1$；而一个标准 W 玻色子有 $\mathfrak{h}_s= (\half;0)$，因而给出 $\Omega(u,\gamma;y)=y+1/y$。事实上，在我们所知的一切具有良好紫外极限的 $\CN=2$ 场论例子中，BPS 粒子所在的表示在 $su(2)_R$ 下的分解都只含单态。因此我们把 $\mathfrak{h}_s$ 含 $su(2)_R$ 非平凡表示的 BPS 粒子称为\emph{奇特 BPS 粒子}（exotic BPS particle）。类似地，我们把在 $su(2)_R$ 下非平凡变换的带框 BPS 态称为\emph{奇特带框 BPS 态}。对非奇特的 BPS 态，我们的 PSC 与 \cite{Gukov:2004hz,Dimofte:2009bv,Dimofte:2009tm} 的「精细化指标」相同，至多差一个（依赖约定的）符号。

现在我们对「原味」与带框 BPS 态各陈述一组「正性猜想」。其中最强的一条就是上面提到的「无奇特粒子猜想」：$\mathfrak{h}_s$ 的同型分解只含 $su(2)_R$ 单态。也就是说，在库仑分支的光滑点上没有奇特粒子。\footnote{在希格斯分支或混合分支上，陈述这样的命题会有一些麻烦。}注意，借由穿墙，奇特 BPS 粒子的存在必然会带来奇特带框 BPS 态。一条较弱的猜想是「强正性猜想」，它断言 $\Omega(u,\gamma;y)$（以及其带框对应物）是 $su(2)$ 特征标 $\chi_n(y)$ 的线性组合，且系数为\emph{非负整数}。「弱正性猜想」则只要求 $\Omega(u,\gamma;y)$ 在 $y=+1$ 处为正整数。

无奇特粒子猜想蕴涵强正性猜想，但反过来未必成立。强正性猜想的一个充分条件是：$\mathfrak{h}_s$ 的所有同型分量都具有整值的 $su(2)_R$ 自旋。然而这个条件并非必要。

我们将在第 4 节% TODO 合稿改 \ref{sec:Formal-line-operators}
中看到，强正性猜想（哪怕只取弱化形式）对带框与标准 BPS 谱都有深远而漂亮的推论。有很强的证据表明，它对 $A_1$ 类中的一切理论都以一种非常非平凡的方式成立。另一个在我们所有例子中似乎都成立的事实是：对固定的线算子，空间 $\CH_{u,L,\zeta,\gamma}^{\rm BPS}$ 只在有限多个电荷 $\gamma$ 上非零。也就是说，对固定的线算子，只有有限多个非零的带框 BPS 不变量。我们没有简单的场论论证来证明这一论断，但我们猜想这个有限性也许是场论中普遍成立的。不过必须承认，我们的全部证据都基于 $A_1$ 类理论。这也是为什么更好地理解 class $\CS$ 中更高阶的理论很重要，我们希望日后回到这一课题。\footnote{最近（2012 年 8 月），Diaconescu 等人证明了：对纯 $SU(K)$ 规范理论，无奇特粒子猜想确实成立。}

最后还有一点观察。自旋特征标近来在文献中已有研究，但总是在场论的语境里，而非超引力的语境里。我们的构造提示了场论情形更受青睐的一个原因：在超引力中没有可用的 $SU(2)_R$ 对称性，因而无法构造保护自旋特征标。

\subsection{halo 与 halo Fock 空间}
\label{subsec:Halos}

现在我们来描述一类特别有趣的带框 BPS 态，即「芯」（core）与包围它的 BPS 粒子「halo」之间的束缚态；这个芯支撑在线算子附近的一个小邻域内。这类构型与 Denef 的 $\CN=2$ 超引力 halo 解密切相关 \cite{Denef:2000nb,Denef:2000ar,Denef:2002ru}。在 \cite{Denef:2007vg} 中，这一现象被用来推导半本原穿墙公式。我们马上会看到，halo 对穿墙问题的带框版本同样重要。

我们先作经典讨论，取最简单的情形：一个电荷为 $\gamma_c$ 的 BPS 芯，束缚着单个电荷为 $\gamma_h$ 的 halo 粒子。\cite{Denef:2000nb} 中一个论证的变体给出如下结论。\footnote{我们在附录 D 中给出推导的若干细节。% TODO 合稿改 \ref{app:Low-Energy-FxdPoint}
}探测与芯电荷 $\gamma_c$ 相联系的背景红外的 halo 粒子，其能量为
\begin{equation}\label{eq:Halo-Particle-Energy}
E_{halo} = \vert Z_{\gamma_h}(u(r))\vert \left(1 + \cos(\alpha_h -
\alpha_\zeta)\right) -  {\rm Re}(Z_{\gamma_h}(u)/\zeta).
\end{equation}
回想 $u \in \CB$ 代表无穷远处真空模的值。探针粒子在与线算子相距 $r$ 处所探测到的模依赖于 $r$，我们记 $Z_{\gamma_h}(u(r)) =\vert Z_{\gamma_h}(u(r))\vert e^{i\alpha_h(r)}$。另一方面，$\zeta$ 与 $r$ 无关，我们把它的相位写作 $\zeta = e^{i \alpha_\zeta}$。

halo 粒子在某个 $r$ 处使能量取极小，那里 $e^{i\alpha_h(r)} / \zeta = -1$。这给出 Denef 关于 BPS 黑洞束缚态半径的公式 \cite{Denef:2000nb} 的一个密切类比：\footnote{的确，用式 (3.8)% TODO 合稿改 \eqref{eq:heuristic-bound}
附近所解释的启发式论证，取中心荷 $\zeta M - Z_{\gamma_c}$、电荷 $\gamma_f + \gamma_c$，Denef 的公式在 $M\to +\infty$ 极限下就化为 \eqref{eq:halo-rad}。}
\begin{equation}\label{eq:halo-rad}
r_{\rm halo} = \frac{\langle \gamma_h, \gamma_c\rangle}{2 {\rm
Im}(Z_{\gamma_h}(u)/\zeta)}.
\end{equation}
这个结果有几个重要的推论。第一，它给出 Denef 稳定性条件的一个类比：halo 构型只对使 \eqref{eq:halo-rad} 给出的 $r_{\rm halo}$ 为正的那些 $(u,\zeta)$ 存在。第二，注意当 $Z_{\gamma_h}(u) = - \vert Z_{\gamma_h}(u)\vert \zeta$ 时 $r_{\rm halo} \to \infty$——换句话说，当 $(u, \zeta)$ 落在 BPS 墙 $\widehat W(\gamma_h)$ 上时。除非我们非常靠近这道墙，$r_{\rm halo}$ 都远小于红外描述成立的自然截断。所以图像是这样的：当人们趋近 BPS 墙 $\widehat W(\gamma_h)$ 时，由（某个倍数）电荷 $\gamma_h$ 的粒子搭成的 halo 构型会达到一个尺度，使得把它们完全放在红外理论里处理是合理的。当我们到达墙上时，这些 halo 胀到无穷大，并在另一侧消失。

与这些 halo 相联系的\emph{量子}态是什么？幸运的是，经典 halo 构型很容易量子化：由于彼此互为 BPS，halo 粒子之间不发生相互作用，它们只是生成一个 $\IZ_2$ 分次的量子态 Fock 空间。为在数学上描述它，设 $\CH'(\gamma_h;u)$ 是电荷为 $\gamma$ 的「原味」BPS 态空间，并已抽去一个半超多重态。它是空间 $so(3)$ 与 $su(2)_R$ 的一个表示。我们用一个 $\IZ_2$ 分次，其符号由 $-(-1)^{2J_3}$ 给出。通常的 BPS 简并度 $\Omega(\gamma_h;u)$ 就是 $\CH'(\gamma_h;u)$ 关于这个分次的超维数取负。令 $J_{\gamma_c,\gamma_h}= \half (\vert \langle \gamma_c, \gamma_h\rangle \vert-1)$，引入空间 $so(3)$ 的表示 $(J_{\gamma_c,\gamma_h})$，其维数为 $\vert \langle \gamma_c, \gamma_h\rangle \vert$。这解释了 halo–芯相互作用所产生的电磁场的自旋。我们在有限维向量空间 $(J_{\gamma_c,\gamma_h})\otimes \CH'(\gamma_h;u)$ 上构造一个 $\IZ_2$ 分次 Fock 空间，其中把 $(J_{\gamma_c,\gamma_h})$ 的 $\IZ_2$ 分次取为\emph{偶}。\footnote{给定一个有限维 $\IZ_2$ 分次向量空间 $V= V^0 \oplus V^1$，超维数为 $(n_0\vert n_1)$，与之相伴的 $\IZ_2$ 分次 Fock 空间是 $V^0$ 上的对称代数与 $V^1$ 上的反对称代数之张量积。更通俗地说，我们取一组基 $\alpha_i$，$i=1,\dots, n_0$，以及 $\gamma_s$，$s=1, \dots, n_1$，然后把 $\alpha_i$ 对应到产生算子的玻色型海森堡代数、把 $\gamma_s$ 对应到产生算子的费米型克利福德（Clifford）代数表示出来。这些看起来古怪的 $\IZ_2$ 分次有一个漂亮的物理解释 \cite{Denef:2002ru}，归根结底与磁场迫使 halo 粒子的自旋朝内这一事实有关。}

具体地说，若我们由
\begin{equation}
\Omega(u,\gamma_h;-z) =  {\rm Tr}_{\CH'(\gamma_h;u)}(-z)^{2J_3} z^{2 I_3} =
\sum_{-M_{\gamma_h}}^{M_{\gamma_h}} a_{m, \gamma_h} z^m
\end{equation}
定义一组整数 $a_{m, \gamma_h}$，其中 $M_{\gamma_h} \ge 0$ 是 halo 粒子的最大 $\CJ_3$ 的两倍，则我们的 Fock 空间由 $\vert a_{m,\gamma_h}\vert$ 个产生算子生成，这些产生算子的 $2 \CJ_3$ 本征值为 $m+m'$，其中 $m'$ 取遍形如 $m' = -2J_{\gamma_c,\gamma_h}, -2J_{\gamma_c,\gamma_h} + 2, \dots,2J_{\gamma_c,\gamma_h}-2, 2J_{\gamma_c,\gamma_h}$ 的值。当 $m$ 为偶数（即 $a_{m,\gamma_h}>0$）时振子是费米型的，当 $m$ 为奇数（即 $a_{m,\gamma_h}<0$）时振子是玻色型的。注意 $\Omega(\gamma_h;u) = \sum_m a_{m, \gamma_h}$。若 $\langle\gamma_c,\gamma_h \rangle=0$，则我们把 $(J_{\gamma_c,\gamma_h})$ 视为零向量空间，此时不形成任何 halo。当然，$a_{m, \gamma_h}$ 是 $u$ 的分段连续整值函数，但我们通常在记号中略去这一依赖。
